\section{Méthodologie}
\label{sec:methodo}

\subsection{Source et description des données}
Les données proviennent de la plateforme \textbf{ENTSO-E Transparency Platform}
\cite{entsoe}, diffusées via le portail open-data \emph{Open Power System Data}
(fichier brut \texttt{data/raw/time\_series\_60min\_singleindex.csv}). Après
sélection des colonnes allemandes, conversion des horodatages en
\texttt{DatetimeIndex} UTC, tri chronologique et dédoublonnage, le jeu
analytique final est \texttt{data/processed/germany\_energy\_analysis.csv}.

\begin{table}[H]
\centering
\caption{Fiche synthétique du jeu de données.}
\label{tab:donnees}
\begin{tabular}{ll}
\toprule
Champ & Valeur \\
\midrule
Thématique & Énergie / Environnement --- Allemagne \\
Source & ENTSO-E Transparency Platform \\
Fréquence & Horaire (60 min), horodatage UTC \\
Période & 2015-01-01 00:00 $\rightarrow$ 2020-09-30 23:00 \\
Nombre d'observations & \(50\,400\) \\
Trous / doublons & \(0\) \\
Valeurs manquantes & \texttt{Load} : 0 ; \texttt{Wind} : 74 ; \texttt{Solar} : 103 \\
Variables retenues & \texttt{Load}, \texttt{Wind}, \texttt{Solar} (MW) \\
\bottomrule
\end{tabular}
\end{table}

\paragraph{Variables retenues.}
\texttt{Load} est la consommation électrique nationale (variable endogène
principale) ; \texttt{Wind} et \texttt{Solar} sont les productions éolienne et
photovoltaïque. Le triplet consommation + deux renouvelables constitue un
système multidimensionnel cohérent, justifiant le volet VAR/VECM.

\paragraph{Contrôles qualité.}
L'index horaire attendu (\texttt{date\_range} de fréquence horaire) coïncide
exactement avec les horodatages présents : aucun trou, aucun doublon. Les
valeurs manquantes sont rares et isolées (74 pour \texttt{Wind}, 103 pour
\texttt{Solar}) ; elles sont \textbf{conservées} dans le fichier analytique
afin de ne pas fausser l'exploration, puis imputées juste avant la
modélisation.

\subsection{Prétraitement et imputation}
La construction de l'objet série temporelle natif diffère légèrement entre R
et Python : côté R, on utilise \texttt{ts(..., frequency = 24)} (cycle
quotidien) ; côté Python, un \texttt{DataFrame} pandas à \texttt{DatetimeIndex}
UTC avec fréquence déclarée \texttt{h}. L'imputation suit une approche
heure-saisonnière identique entre les deux langages : médiane par heure de la
journée, complétée par une interpolation linéaire
(\texttt{forecast::na.interp} en R). Ce choix préserve la régularité horaire
indispensable aux modèles saisonniers (\(m=24\)), contrairement à une
suppression de lignes qui briserait le pas de temps.

\subsection{Découpage chronologique}
Le découpage est \textbf{strictement chronologique} \(90\,\%\) / \(10\,\%\)
(jamais aléatoire, pour éviter toute fuite d'information future) :
\(45\,360\) heures d'apprentissage et \(5\,040\) heures de test. Le jeu de test
débute le \textbf{5 mars 2020} et inclut donc l'entrée dans la crise sanitaire
COVID-19, ce qui constitue un point d'interprétation important. Quatre
horizons sont évalués : \(h \in \{5,\,10,\,24,\,5\,040\}\) heures (\(5\) et
\(10\) pour le court terme, \(24\) pour le cycle journalier, \(5\,040\) pour la
trajectoire complète sur tout le jeu de test).

\subsection{Analyse exploratoire (EDA)}
L'EDA comprend : (i) l'évolution temporelle des trois séries sur toute la
période ; (ii) les profils moyens horaire, hebdomadaire, mensuel et annuel ;
(iii) la cartographie des valeurs manquantes ; (iv) la détection d'observations
atypiques (règle IQR globale et MAD roulante sur \(24\)~h) ; (v) les
corrélogrammes ACF/PACF (séries brutes et différenciées) ; (vi) une
décomposition STL (tendance / saisonnalité / résidus) pour la charge.

\subsection{Stationnarité et ordres d'intégration}
La stationnarité est testée par le test \textbf{ADF} (H0 : racine unitaire) et
le test \textbf{KPSS} (H0 : stationnarité), sur les séries brutes, en
différence première (\(d=1\)) et en différence saisonnière (\(d=24\)). La
conclusion retenue exige l'accord des deux tests.

\subsection{Modèles univariés : ARIMA et SARIMA}
Pour chaque série, trois candidats sont comparés par critère d'information
(AIC/AICc/BIC) :
\begin{enumerate}
  \item \textbf{SARIMA\_manuel} : \(\text{ARIMA}(1,1,1)(1,1,1)_{24}\), motivé
        par les corrélogrammes de l'EDA (rebonds à 24 et 168) ;
  \item \textbf{ARIMA\_auto} : \texttt{auto.arima} non saisonnier (\(d=1\)) ;
  \item \textbf{SARIMA\_auto} : \texttt{auto.arima} saisonnier borné à
        \(m=24\).
\end{enumerate}
Le modèle retenu est celui qui minimise l'\textbf{AICc}. Les résidus sont
diagnostiqués par le test de \textbf{Ljung--Box} (autocorrélation, retard 24),
le test de \textbf{Jarque--Bera} (normalité) et un test \textbf{ARCH--LM}
(12 retards) implémenté manuellement. Les prévisions sont produites aux
horizons \(h=5\), \(10\), \(24\) et \(5\,040\) avec les métriques MSE, RMSE,
MAE et MAPE (ce dernier calculé sans division par zéro, \texttt{Solar} étant
nulle la nuit).

\paragraph{Intervalles de prévision.}
Pour le modèle retenu, les \textbf{intervalles de prévision} à \(95\,\%\)
sont calculés aux horizons \(h=5,\,10,\,24\) (c'est-à-dire dans les plages où
la spécification reste valide, contrairement à \(h=5\,040\) où la bande
diverge). Ils sont obtenus via \texttt{pmdarima} pour les spécifications
saisonnières : le mode \texttt{low\_memory} de \texttt{statsmodels},
nécessaire pour ajuster un SARIMA saisonnier sur \(45\,360\) points sans
saturer la mémoire, n'expose pas la covariance de prévision.

\subsection{Modèles multivariés : VAR, Granger, Johansen, VECM}
Les liaisons structurelles entre consommation et productions s'expriment à
l'échelle \textbf{journalière} : les séries sont donc agrégées en moyennes
quotidiennes (\(\approx 2\,100\) jours). Pour chaque série univariée, le test
ADF justifie la différenciation \(D1\), puis :
\begin{itemize}
  \item l'ordre \(p\) du \textbf{VAR} est sélectionné par critère d'information
        (AIC/HQ/SC/FPE, \texttt{lag.max}=15) ;
  \item la \textbf{causalité de Granger} est testée en paires pour les six
        relations ordonnées (H0 : \(X\) ne cause pas \(Y\) au sens de Granger) ;
  \item la \textbf{cointégration de Johansen} (test de trace) est menée sur
        les niveaux pour déterminer le rang \(r\) ; si \(r\geq 1\), un
        \textbf{VECM} est estimé (vecteurs \(\beta\) normalisés et vitesses
        d'ajustement \(\alpha\)) ;
  \item les dynamiques sont synthétisées par les \textbf{fonctions de réponse
        impulsionnelle} (IRF, choc \texttt{Wind}/\texttt{Solar} sur
        \texttt{Load}, 28 jours) et la \textbf{décomposition de la variance}
        (FEVD).
\end{itemize}

\subsection{Modélisation de la volatilité : ARCH / GARCH}
Les tests ARCH--LM ayant rejeté l'absence d'effet ARCH, la variance
conditionnelle des variations journalières est modélisée par trois
spécifications comparées, avec innovations de Student (queues épaisses) :
\textbf{GARCH(1,1)}, \textbf{GARCH(2,1)} et \textbf{EGARCH(1,1)}. La
\textbf{persistance} est définie comme la somme de tous les coefficients ARCH
et GARCH pour un sGARCH (respectivement \(\beta_1\) pour un eGARCH) ; la
demi-vie de la volatilité est \(\log(0{,}5)/\log(\text{persistance})\)
lorsque la persistance est strictement comprise entre 0 et 1.

\subsection{Modèles profonds : LSTM et GRU}
En Python (Keras/TensorFlow), des réseaux récurrents \textbf{LSTM} et
\textbf{GRU} sont entraînés séparément avec : fenêtre glissante de
\(168\)~h (\(=1\) semaine), normalisation \emph{MinMax ajustée sur le seul
train}, validation chronologique de \(5\,040\)~h, \textbf{couche de Dropout}
(\(p=0{,}2\)) entre l'état récurrent et la sortie pour limiter le
surapprentissage, et prévision récursive multi-pas. Les deux architectures
sont évaluées et la meilleure (sur la \texttt{val\_loss}) est conservée par
série ; l'évaluation est menée à un pas en \emph{walk-forward} et sous forme
de trajectoires aux horizons \(h=5/10/24/5040\).

\paragraph{Validation croisée temporelle (TSCV).}
En complément du simple hold-out, une \textbf{validation croisée à fenêtre
expansive} (\(k=3\) plis) est appliquée aux réseaux : le train grandit
progressivement, la validation reste toujours postérieure au train, ce qui
fournit une estimation de la RMSE plus robuste et sans fuite temporelle.

\subsection{Intervalles de prévision}
Les intervalles à \(95\,\%\) sont produits pour le modèle retenu à
\(h=5,\,10,\,24\) (SARIMA via \texttt{pmdarima}) et pour Prophet
(nativement, bande \texttt{yhat\_lower}/\texttt{yhat\_upper}). Les modèles
profonds et l'hybride fournissent une prévision ponctuelle ; l'incertitude y
est appréhendée par la dispersion des erreurs walk-forward.

\subsection{Modèles additionnels : Prophet et hybride SARIMA+LSTM}
Deux modèles complémentaires confrontent les approches classiques et
profondes :
\begin{itemize}
  \item \textbf{Prophet} (Meta) : décomposition automatique en tendance
        linéaire par morceaux, saisonnalités quotidienne / hebdomadaire /
        annuelle et jours fériés, robuste aux valeurs manquantes ;
  \item \textbf{hybride SARIMA\,$+$\,LSTM} : la composante \emph{linéaire}
        est captée par le SARIMA retenu, dont les résidus (supposés porter la
        structure non linéaire) alimentent un LSTM ; la prévision hybride est
        la somme des deux, méthode classique d'hybridation
        tendance\,/\,résidu.
\end{itemize}


\subsection{Évaluation comparative}
Les métriques sont MSE, RMSE, MAE et MAPE. Deux protocoles sont distingués :
(i) des \textbf{trajectoires} émises depuis une origine unique aux horizons
\(h=5/10/24/5040\) ; (ii) un \textbf{walk-forward 1 pas} sur les \(5\,040\)~h
de test (origine glissante), où chaque prévision n'utilise que l'information
passée. Le walk-forward confronte les baselines \emph{naïve} et
\emph{saisonnière-naïve} (\(24\)~h) au SARIMA retenu, au LSTM et au GRU.

Les comparaisons par paires utilisent le test de \textbf{Diebold--Mariano}
sur les erreurs quadratiques, appliqué à \textbf{toutes les paires} de
modèles d'une même série, avec correction \textbf{HLN}
(Harvey--Leybourne--Newbold) et variance \textbf{HAC} (Newey--West, retard
\(h-1\)). Enfin, un \textbf{tableau de comparaison} enrichi associe aux
performances la \emph{complexité} (nombre de paramètres), le \emph{coût de
calcul} et l'\emph{interprétabilité} de chaque famille de modèles, comme
demandé au cahier des charges.
